\section{为什么要研究模型的“生物学”}

这场讲座的标题不是修辞装饰，而是一种研究立场：工程师写下 Transformer 架构、训练目标和优化流程，却没有逐行写出模型最终学会的内部算法。训练像演化一样产生复杂对象，研究者面对的是一个“长出来”的系统。于是 mechanistic interpretability（机制可解释性）的任务，不是给输出贴标签，而是把一次前向传播拆成可检验的部件、连接和因果假设。

\lecturefigure{slide-01-title-biology-large-language-model.jpg}{讲座标题：把训练后的大模型当作复杂系统来解剖}{00:01:24}

\teachervoice{00:01:41--00:02:43，老师把 interpretability 与 biology 类比：两者都研究由某种优化或演化过程产生的复杂对象。这个类比不是说神经网络是生命，而是提醒我们：内部机制并非按人类模块图逐项装配。}

\subsection{能力惊喜与机制怪异必须同时解释}

如果只看成功案例，模型像一个能在上下文中临时学会新任务的通用程序；如果只看失败案例，它又像一个会把互相矛盾的事实拼在一起的随机系统。课堂先故意把这两面并置，因为真正的研究问题不是“模型聪明还是愚蠢”，而是同一套参数如何在不同输入上组织出截然不同的计算路径。

\lecturefigure{slide-02-models-do-cool-things.jpg}{低资源 Circassian 翻译：上下文示例足以诱导新的语言能力}{00:04:04}

图中案例把多年积累的俄语--Circassian 对照表直接放入上下文，模型不仅完成新句翻译，还能分析语法。应先比较“专门训练模型”与“通用模型临时读取词表”这两条路线：后者没有重新更新参数，却在一次上下文中组合已有语言表示。这支持强大的抽象迁移能力，但不证明模型已经可靠掌握该语言的全部文化、语法和事实。

\lecturefigure{slide-03-models-do-weird-things.jpg}{闰日矛盾：正确事实、错误前提与局部推理被同时激活}{00:05:20}

这张图的关键不是一个日期错误，而是输出内部的冲突结构：模型知道 2024 是闰年规则的实例，又生成“2024 不是闰年”，随后还正确说下一个闰年是 2028。局部事实检索与整体一致性并未自动合并成单一算法。它提示我们，流畅文本可能是多个部分策略竞争后的表面结果。

\begin{warningbox}{能力提升不会自动消灭怪异}
老师强调，模型变强以后，明显的“六根手指”式错误会减少，但怪异往往只是变得更隐蔽、更高层。Benchmark 提升只能说明某些行为更好，不能说明内部冲突已经被理解或消除。
\end{warningbox}

\lecturefigure{slide-04-three-core-findings.jpg}{本讲三条主线：抽象表示、并行计算与规划}{00:08:19}

读这张路线图时，应把左侧常见说法与右侧研究发现逐行对照：模型并非只匹配训练样本，而会组合抽象表示；并非只执行浅层串行规则，而能在一次前向传播中并行运行多条路径；并非只能响应当前 token，也会形成影响未来词形和结构的计划。后续每个案例都要用 intervention（干预）而非相关性截图来支撑这些主张。

这三条发现还对应三种不同的错误直觉。把模型理解为检索器，会漏掉 feature 之间的组合；把模型理解为逐步执行自然语言程序，会漏掉同一层内同时存在的候选路径；把 next-token prediction 理解为“只看下一词”，会漏掉隐藏状态中对更远目标的约束。讲义后面不会把任何单个案例升级成普遍定律，而是反复检查：该图来自哪个 prompt，替换模型解释了多少，干预改变了什么，以及是否存在同样合理的替代回路。

\begin{importantbox}{本讲的证据问题}
一个内部解释至少要回答三件事：\textbf{表示了什么}、\textbf{如何连接成计算}、\textbf{改变该部件是否会改变行为}。只有 feature label 没有回路，是字典；只有相关图没有干预，是假说；干预成功但覆盖不完整，也仍不是全局解释。
\end{importantbox}

\subsection{本章小结}

“生物学”视角把模型从 API 黑盒变成待解剖的复杂系统。开场案例建立了整份讲义的张力：大模型能快速组合新能力，也会把正确与错误策略同时带入答案；因此需要从行为评测继续深入到表示、回路和因果验证。

